% ============================================================================
\chapter{Introduction}\label{ch:intro}
% ============================================================================

\section{The reconstruction problem}

Every quantitative desk begins its day with the same reconstruction problem. The market does not
publish an implied volatility \emph{surface}; it publishes a few thousand option quotes sparse in
some regions, dense in others, noisy everywhere, and living on an irregular grid of strikes and
expiries that changes daily. Everything downstream of those quotes, including marking books, computing Greeks, calibrating stochastic volatility and local volatility models, and stress testing portfolios, relies not on the quotes themselves but on a \emph{smoothed and completed surface} constructed from them.
 The construction of that
surface is therefore not a preprocessing detail: it is the first model in the pricing chain, and
its errors propagate into every model behind it \cite{fengler2009}.

The construction is constrained. A surface that fits the quotes but admits \emph{static arbitrage}, such as a butterfly with a nonnegative payoff but a negative price or a calendar spread exhibiting the same defect, is not merely inelegant. By the Breeden--Litzenberger identity
\cite{breeden1978} it corresponds to a negative risk-neutral density; by Dupire's formula
\cite{dupire1994} it corresponds to a negative local variance; operationally it produces false
Greeks and phantom P\&L \cite{fengler2009}. The tension of the field is that the two goals
conflict: fitting the quotes pulls toward flexibility, excluding arbitrage pulls toward rigidity.

Three families of approaches coexist. \emph{Parametric} surfaces, most notably SVI \cite{gatheral2004} and its surface extension SSVI \cite{gatheraljacquier2014}, provide closed-form derivatives and, in the case of SSVI, explicit no-arbitrage conditions, but rely on a relatively rigid functional form. \emph{Per-surface neural smoothers} \cite{ackerer2020,hoshisashi2023} retain a parametric prior and augment it with a neural correction term. No-arbitrage is imposed through \emph{soft penalties} applied to surface derivatives, which are computed exactly using automatic differentiation. \emph{Neural operators} \cite{wiedemann2025,lu2021,kovachki2023} extend this idea by learning the smoothing \emph{map} from quote configurations to complete surfaces, rather than fitting a separate model for each day. This amortises training over the full historical dataset and allows a new surface to be evaluated in milliseconds.


\section{The research gap: what do soft constraints actually guarantee?}

The thesis was designed as a controlled benchmark: the three families implemented on one dataset,
under one protocol, with one audit. That benchmark exists and is reported in full. But building
it surfaced a sharper question. Midway through the derivative-constrained smoother, the training
log displayed two numbers side by side. The first was \texttt{pen\_cal}, the calendar-arbitrage
term of the training loss: the average of $\operatorname{ReLU}(-\partial_\tau w_\theta)^2$ over
the calendar collocation nodes, which is positive wherever total variance decreases in maturity
(a calendar violation) and zero wherever it does not. It read \texttt{pen\_cal = 0.000e+00}, meaning that at every node of its collocation grid the fitted surface had a
non-negative maturity slope and the penalty had nothing to punish. The second number came from
an independent audit grid, finer than and disjoint from the collocation grid, and reported a
$28\%$ calendar-violation rate on the very same surface. Both numbers were true: the penalty was
satisfied exactly \emph{where it was sampled}, and the surface violated the constraint \emph{between}
the sampled nodes. The penalty
had been minimised to zero \emph{where it was sampled}, and the surface violated \emph{between
the samples}, because the collocation grid, exponentially spaced in maturity as prescribed for
the butterfly constraint, left a $0.41$-year hole exactly where calendar arbitrage lives.

This observation, which gives the most literal interpretation of Chataigner et al.'s warning that penalties act ``only on the grid'' \cite{chataigner2020}, reshaped the direction of the thesis. Soft constraints are ultimately statements about a particular sampling set. The relevant question for any constrained smoother is therefore not \emph{whether the penalty is zero}, but \emph{where it is zero, how that set relates to the regions in which the constraint may fail, and what such failures represent in economic terms}.\\


The thesis answers this question three times, at increasing levels of generality:
\begin{enumerate}
\item \textbf{Diagnosis and repair} (Section~\ref{sec:deep}): the blind spot is exhibited,
      explained by an interaction between the geometry of the two arbitrage constraints and the
      grid that serves them, and repaired by \emph{constraint-specific collocation}. On a stress
      test whose own validity is verified by a two-condition gate, material calendar violations
      fall from $34.4\%$ to $0.0\%$.
\item \textbf{Transversality} (Section~\ref{sec:operator}): the neural operator inherits the
      pathology in a stronger form---blind between nodes \emph{and} across days. Its own
      penalty-grid audit reports near-zero violations on $386$ out-of-sample days; an independent
      fine audit flags $330$ of them as materially butterfly-violating.
\item \textbf{Pricing the failure} (Chapter~\ref{ch:downstream}): through
      $\sigma_{\mathrm{loc}}^2=\partial_\tau w/g$, the audited quantities \emph{are} the pricing
      quantities. Violations become negative local variances requiring repair on up to $50\%$ of
      the operator's surface, negative density mass up to $230\%$, and a measurable degradation
      of residual-signal quality.
\end{enumerate}

\paragraph{Constructive counterpart.} The critique has a constructive twin. The
Gatheral--Jacquier theory is used not only as an audit but as a \emph{certificate}: our SSVI
prior is calibrated \emph{inside} the no-arbitrage region of \cite{gatheraljacquier2014}
(Theorem~\ref{lem:range}), so the ablation finding ``the prior does the protective work'' rests on a
theorem, not on luck. Symmetrically, every claim about a stress test or an invariance property is
\emph{gated}: a fit-validity gate certifies that a stressor was actually fitted before its
arbitrage comparison is read, and a collapsed operator's trivially perfect discretization
invariance is excluded rather than reported. The theoretical chapter adds two elementary but,
to our knowledge, previously unstated formalisations of the mechanism: a \emph{node-gap bound}
(Theorem~\ref{prop:nodegap}) showing that the depth of a calendar violation between
collocation nodes scales \emph{linearly} with the node gap, and a curvature-inertness argument showing that
piecewise-affine correctors make the butterfly penalty structurally blind to the very derivative
it constrains (Section~\ref{sec:softtheory}).

\section{Contributions}

\begin{enumerate}[label=(\roman*)]
\item A validated replication of the Gatheral--Jacquier no-arbitrage theory and calibration
recipes, matched to the paper's own numbers to $10^{-7}$, including the identification and
correction of a missing term in the published inverse of the natural-parameter map
(Section~\ref{sec:parametric}, Remark~\ref{rem:gjtypo});
\item a derivative-constrained deep smoother with exact-autodiff penalties, and the diagnosis,
mechanism and repair of the \emph{collocation blind spot}, including the node-vs-audit gap and a
materiality threshold as standard diagnostics (Section~\ref{sec:deep});
\item a certified-arbitrage-free SSVI prior via an endpoint-monotonicity argument
(Theorem~\ref{lem:range}), together with the node-gap bound (Theorem~\ref{prop:nodegap})
formalising why one collocation grid cannot serve two constraints;
\item an executable operator-smoothing study (DeepONet-class and graph-neural-operator
architectures) with a chronological 1{,}926-day design, a synthetic-pretraining transfer
experiment, and the demonstration that the blind spot is transversal and two-dimensional
(Section~\ref{sec:operator});
\item a downstream-economics layer that prices the audit: Dupire local volatility,
Breeden--Litzenberger densities, Monte-Carlo repricing, a data-level static-arbitrage scanner on
bid/ask quotes, and a residual relative-value backtest with an arbitrage-dirty counterfactual
(Chapter~\ref{ch:downstream}).
\end{enumerate}

\section{Related work and positioning}\label{sec:related}

\paragraph{Parametric surfaces.} SVI \cite{gatheral2004} owes its longevity to three
properties: closed-form derivatives, trader-interpretable parameterisations (raw, natural and
Jump--Wings), and exact large-maturity consistency with the Heston model. Gatheral and Jacquier
\cite{gatheraljacquier2014} supply the modern theory, and it is the backbone of
Chapter~\ref{ch:theory}: explicit butterfly conditions expressed through the Durrleman function
$g$, exact conversions between the three parameterisations, an algorithm eliminating butterfly
arbitrage in Jump--Wings space, the SSVI surface with closed-form static-arbitrage conditions
(Theorems~\ref{thm:gjcal} and~\ref{thm:gjbfly}), and arbitrage-free interpolation and
extrapolation recipes. Section~\ref{sec:parametric} replicates all of it and validates each
component against the numbers reported in the original paper.

\paragraph{Arbitrage-free smoothing before machine learning.} The problem of producing a surface
that is simultaneously faithful to the quotes and free of static arbitrage predates any learned
model, and the early literature solved it by choosing model classes on which the constraints are
tractable. Kahale \cite{kahale2004} interpolates directly in price space using
convexity-preserving splines, so that the butterfly condition cannot be violated by
construction. Fengler \cite{fengler2009} formulates smoothing as a constrained optimisation over
natural cubic splines with the no-arbitrage inequalities imposed as hard linear constraints, and
documents the operational cost of ignoring them: negative transition probabilities and
misleading Greeks. Cont and da Fonseca \cite{cont2002} characterise the empirical dynamics of
the surface itself, a reminder that a smoother must be stable from one day to the next and not
merely accurate on a single cross-section, which is why day-over-day stability appears among the
downstream metrics of Chapter~\ref{ch:downstream}. Two features of this literature carry
through: the constraints are imposed \emph{exactly}, and the model class is chosen so that the
constraint set remains tractable. The learned smoothers below relax both, and that is precisely
where this thesis intervenes.

\paragraph{Neural networks for option surfaces: a short history.} Learned option pricing begins
with Hutchinson, Lo and Poggio \cite{hutchinson1994}, who fit pricing and hedging functions
nonparametrically and already observe that an unconstrained network can reproduce prices while
violating the qualitative shape properties theory requires. Dugas et al.\ \cite{dugas2001} give
the first influential answer, building the constraints into the \emph{architecture} through
functional forms that are monotone and convex by construction. The alternative answer, imposing
constraints through the \emph{loss} instead of the architecture, arrives from the numerical-PDE
literature. Lagaris et al.\ \cite{lagaris1998} train networks to satisfy differential equations
by penalising the residual at sampled points, and physics-informed neural networks
\cite{raissi2019} turn this into a general methodology resting on one observation: because a
network's derivatives are available \emph{exactly} through automatic differentiation
\cite{baydin2018}, a differential constraint can be evaluated pointwise and penalised during
training. That observation is what makes derivative-constrained smoothing possible at all, and
every soft-constrained model in this thesis inherits it. It also carries the limitation this
thesis measures, and the limitation is visible already in the PINN formulation: the residual is
penalised on a finite set of collocation points, so what training enforces is a statement about
that set.

\paragraph{Per-surface neural smoothers.} Ackerer et al.\ \cite{ackerer2020} bring the
prior-times-corrector design to implied volatility: a parametric surface multiplied by a neural
correction, with the no-arbitrage conditions imposed as soft penalties on an extended
collocation grid, their constraint set $\mathcal I_{C4,5,6}$ covering calendar arbitrage,
butterfly arbitrage and far-wing linearity. Their sweep over $\lambda\in\{0,1,10\}$ quantifies
the fit-versus-arbitrage trade-off that Section~\ref{sec:deep} reproduces and extends, and their
architecture is the one this thesis certifies and then transplants into an operator. Hoshisashi
et al.\ \cite{hoshisashi2023} supply the training mechanics adopted here, exact nested
derivatives by automatic differentiation and penalties evaluated on a mesh distinct from the
quotes, together with the candid acknowledgement that soft constraints do not fully secure the
no-arbitrage conditions. Chataigner et al.\ \cite{chataigner2020} work on the Dupire side and
state the limitation this thesis makes precise: the penalties act only on the sampled grid.
Their closing remark that extraction \emph{from implied volatilities} is left to a further paper
identifies the gap addressed in Chapter~\ref{ch:downstream}.

\paragraph{From networks to operators.} A network learns one function; an operator learns a map
between function spaces, so a single trained model applies to a whole family of inputs. Chen and
Chen \cite{chenchen1995} prove the universal-approximation theorem for nonlinear operators that
underwrites the programme, and DeepONet \cite{lu2021} turns it into a practical architecture by
separating the encoding of the input function (branch) from the query coordinates (trunk). In
parallel, Li et al.\ represent the operator as a composition of learned kernel integrals,
discretised either by message passing over a graph \cite{li2020gno} or in Fourier space
\cite{li2021fno}; Kovachki et al.\ \cite{kovachki2023} unify these constructions and formalise
discretization invariance. Operator Deep Smoothing \cite{wiedemann2025} applies the
graph-neural-operator variant to implied volatility, mapping a day's quote cloud directly to a
smoothed surface, trained on roughly a decade of proprietary data, evaluating in milliseconds,
with soft no-arbitrage penalties resolved by finite differences on a grid.
Section~\ref{sec:operator} reproduces that design in the academic-data regime, adds a
synthetic-pretraining transfer axis, and audits the penalties on an independent grid.

\paragraph{Downstream and economic layers.} Breeden and Litzenberger \cite{breeden1978}, Dupire
\cite{dupire1994} and Gatheral \cite{gatheral2011book} provide the identities developed in
Chapter~\ref{ch:theory}; Durrleman \cite{durrleman2010} and Roper \cite{roper2010}, together
with Cox and Hobson \cite{coxhobson2005} and Kellerer \cite{kellerer1972}, supply the
characterisation results on which the audit rests. Lee's moment formula \cite{lee2004}
constrains the admissible behaviour of the wings. More recent work connects surface construction
to trading directly: Duan \cite{duan2026} builds a delta-neutral strategy on residuals from a
parametric model augmented by a neural network, and the recent SANOS framework \cite{sanos2026}
proposes a non-parametric arbitrage-free smoother on the same OptionMetrics universe. Together
these indicate that current research is active on precisely the territory studied here.

\paragraph{Positioning.} Relative to this literature, the thesis makes three contributions.
First, it provides a controlled \emph{cross-family} comparison on a common dataset under a
unified evaluation protocol, which none of the cited studies undertakes. Second, to the best of
our knowledge, it is the first study to \emph{measure} the collocation blind spot rather than
caveat it: the gap between what the penalty grid reports and what an independent audit grid
finds, the materiality of the resulting violations, and the frequency of blind days. It further
shows the phenomenon is transversal, appearing both in per-day smoothers and in neural
operators, and in two directions, between grid nodes in space and across days in time.

Third, the thesis insists that a violation be reported in units a practitioner can act on rather
than as a percentage of grid cells. We call this \emph{pricing the audit}, and the bridge is the
Dupire identity of Proposition~\ref{prop:dupire}, $\sigma^2_{\mathrm{loc}}=\partial_\tau w/g$:
the two quantities the audit monitors are exactly the numerator and the denominator of the local
variance a desk simulates with. A butterfly violation is therefore not merely $g<0$ on some
fraction of a grid; it is a region in which the local variance is negative or unbounded, so a
Monte-Carlo scheme cannot be initialised there without overriding the model by hand.
Chapter~\ref{ch:downstream} reports exactly that: the share of the surface requiring such
repair, the negative mass carried by the implied densities, and whether an arbitrage-dirty
surface degrades the information content of the residuals computed against it. This is what
\emph{economic} means in this thesis. It is deliberately distinct from the claim that a
violation is \emph{tradable}, which is a separate question about executable bid-ask quotes and
is treated separately in Section~\ref{sec:scanner}.

\section{Outline}

Chapter~\ref{ch:theory} develops the model-free theory of static arbitrage and the parametric
surfaces built on it. It covers arbitrage conditions in $(k,w)$ coordinates, the Durrleman condition, the implied risk-neutral density, Dupire's formula, the SVI and SSVI parameterisations, the full Gatheral and Jacquier certificates and their proofs, the theory of soft constraints and sampling sets, including the node-gap bound and the materiality calculus, neural-network approximation, and neural operators.

Chapter~\ref{ch:data} presents the data-processing pipeline and the unified evaluation methodology applied to all model families. Chapter~\ref{ch:methods} implements and evaluates the three methodological families. It covers the parametric benchmark in Section~\ref{sec:parametric}, the derivative-constrained deep smoother and the collocation blind spot in Section~\ref{sec:deep}, and the neural operators in Section~\ref{sec:operator}. Chapter~\ref{ch:downstream} evaluates the downstream economic consequences of the audit through local volatility, implied densities, the data-level arbitrage scanner, the residual backtest, and VIX, before bringing the different evaluation layers together. Chapter~\ref{ch:conclusion} concludes. Appendix~\ref{app:repro} links every figure and table to the notebook used to generate it.

